\section{Objetivos}

\subsection{Objetivo General}
Desarrollar un modelo predictivo basado en técnicas de Machine Learning que funcione como un sistema de alertas tempranas, capaz de identificar contratos públicos de infraestructura con alto riesgo de incumplimiento o asociados a procesos irregulares, en el marco del Sistema General de Regalías (SGR). Para ello, se empleará como fuente principal de información la base de datos de la plataforma SECOP II.

\subsection{Objetivos Específicos}
\begin{enumerate}
    \item Identificar y seleccionar variables relevantes de tipo contractual y financiero —como anticipos, modalidad de contratación, número de participantes y ubicación geográfica— que influyan en el riesgo de incumplimientos e irregularidades en contratos de infraestructura pública.

    \item Recolectar, integrar y preprocesar los datos provenientes de la plataforma SECOP II, garantizando su calidad, consistencia y confiabilidad para el análisis predictivo.

    \item Diseñar y entrenar modelos de aprendizaje supervisado, tales como Random Forest, Regresión Logística y XGBoost, con el fin de predecir de manera precisa factores de riesgo en contratos públicos.

    \item Evaluar el desempeño de los modelos mediante técnicas como la validación cruzada estratificada y el uso de métricas robustas, tales como AUC, precisión, recall y F1-score, para determinar su efectividad.

    \item Comparar los resultados obtenidos por los diferentes modelos con el propósito de identificar aquel que ofrezca la mayor capacidad predictiva en la detección temprana de riesgos.

    \item Proponer un esquema de implementación del modelo como herramienta de alertas tempranas, orientado a fortalecer la detección de anomalías por parte de entidades de control y mecanismos de veeduría ciudadana.
\end{enumerate}